% Part: proof-theory
% Chapter: sequent-calculus
% Section: quantifiers

\documentclass[../../../include/open-logic-section]{subfiles}

\begin{document}

\olfileid{pt}{seq}{qua}

\olsection{正则证明与代入}

量词规则由于适用于 \LeftR{\lexists} 和~\RightR{\lforall} 规则的“本征变元条件”而引入了一些复杂性。大部分复杂性可以通过确保这些推理中的常项符号~$a$ 不被重复使用来化解。我们引入以下定义。

\begin{defn}
在 \Log{G1c} 或 \Log{G3c} 中，推导~$\pi$ 是\emph{正则的}，如果
\begin{enumerate}
  \item 没有任何本征变元~$a$ 被用作不止一个 \LeftR{\lexists} 或~\RightR{\lforall} 推理的本征变元，且
  \item 如果 $a$ 是 \LeftR{\lexists} 或~\RightR{\lforall} 推理的本征变元，则它在~$\pi$ 中只出现在该推理上方。
\end{enumerate}
\end{defn}

\begin{lem}\ollabel{lem:subst-regular} 如果 $\pi$ 是正则的，其末尾相继式为
$\Gamma \Sequent \Delta$，$c$ 是在~$\pi$ 中未被用作本征变元的常项符号，且 $t$~是不包含~$\pi$ 的任何本征变元的项，那么将~$\pi$ 中 $c$ 的每次出现替换为~$t$ 所得的 $\Subst{\pi}{t}{c}$ 是一个正则推导。
$\Subst{\pi}{t}{c}$ 的末尾相继式是 $\Subst{\Gamma}{t}{c}
\Sequent \Subst{\Delta}{t}{c}$，其中 $\Subst{\Gamma}{t}{c} =
\Setabs{!A(t)}{!A(c) \in \Gamma}$。
\end{lem}

\begin{proof}
  对 $\pheight{\pi}$ 进行归纳。如果 $\pheight{\pi} = 0$，则 $\pi$ 仅由一条公理组成。那么 $\Subst{\pi}{t}{c}$ 也是一条公理，因为其中的任何公式~$!A(c)$ 都变成了 $!A(t)$。

  如果 $\pheight{\pi} > 0$，我们根据~$\pi$ 的最后一个推理来区分情况。结构规则（在 \Log{G1c} 中）和命题联结词的逻辑规则的情况很简单，留作练习。我们考虑 $\pi$ 以量词推理结尾的情况。
  \begin{enumerate}
    \item 如果最后一个推理是 $\RightR{\lforall}$，则 $\pi$ 具有形式
    \[
    \AxiomC{}
    \RightLabel{$\pi'$}
    \Deduce$\Gamma(c) \fCenter \Delta'(c), !A(a,c)$
    \RightLabel{\RightR{\lforall}}
    \UnaryInf$\Gamma(c) \fCenter \Delta'(c), \lforall[x][!A(x,c)]$
    \DisplayProof
    \]
    根据归纳假设，$\Subst{\pi'}{t}{c}$ 是 $\Gamma(t) \fCenter \Delta'(t), !A(a,t)$ 的一个正则推导。由于 $a$
    是~$\pi$ 的一个本征变元，$a$~不在~$t$ 中出现。
    因此，$\Subst{\pi}{t}{c}$ 中的最后一个推理，
    \[
    \AxiomC{}
    \RightLabel{$\Subst{\pi'}{t}{c}$}
    \Deduce$\Gamma(t) \fCenter \Delta'(t), !A(a,t)$
    \RightLabel{\RightR{\lforall}}
    \UnaryInf$\Gamma(t) \fCenter \Delta'(t), \lforall[x][!A(x,t)]$
    \DisplayProof
    \]
    是合法的：结论不包含~$a$。
    \item 如果最后一个推理是 \Log{G3c} 中的 $\LeftR{\lforall}$，
    则 $\pi$ 具有形式
    \[
    \AxiomC{}
    \RightLabel{$\pi'$}
    \Deduce$!A(s(c),c), \lforall[x][!A(x,c)], \Gamma(c) \fCenter \Delta'(c)$
    \RightLabel{\RightR{\lforall}}
    \UnaryInf$\lforall[x][!A(x,c)], \Gamma(c) \fCenter \Delta'(c)$
    \DisplayProof
    \]
    根据归纳假设，$\Subst{\pi'}{t}{c}$ 是 $!A(s(t),t), \lforall[x][!A(x,t)], \Gamma(t) \fCenter
    \Delta'(t)$ 的一个正则推导。
    $\Subst{\pi}{t}{c}$ 中的最后一个推理，
    \[
    \AxiomC{}
    \RightLabel{$\Subst{\pi'}{t}{c}$}
    \Deduce$!A(s(t),t), \lforall[x][!A(x,t)], \Gamma(t) \fCenter \Delta'(t)$
    \RightLabel{\RightR{\lforall}}
    \UnaryInf$\lforall[x][!A(x,t)], \Gamma(t) \fCenter \Delta'(t)$
    \DisplayProof
    \]
    是合法的。\Log{G1c} 中 $\LeftR{\lforall}$ 的情况类似，
    但稍微简单一些。
    \item 最后一个推理是 $\RightR{\lexists}$ 或 $\LeftR{\lexists}$：留作练习。
  \end{enumerate}
  在每种情况下，我们都在一个正则推导（或在双前提规则下，两个正则推导）的末尾添加了一个相继式。新的相继式与 $\pi$ 的末尾相继式的区别仅在于将 $c$ 替换为了项~$t$。此外，$\pi$ 与 $\Subst{\pi}{t}{c}$ 的本征变元相同。既然假设 $t$ 不包含~$\pi$ 的本征变元，新的末尾相继式也就不包含 $\Subst{\pi}{t}{c}$ 的本征变元。因此，$\Subst{\pi}{t}{c}$ 是正则的。
\end{proof}

\begin{prop}
如果 $\pi$ 是 \Log{G1c} 或 \Log{G3c} 中的一个推导，那么存在一个正则推导~$\pi'$，它具有相同的结构（特别是相同的高度），且仅通过替换本征变元而与 $\pi$ 不同。
\end{prop}

\begin{proof}
仅出于证明的目的，我们称~$\pi$ 中的一个本征变元推理是\emph{干净的}，如果其本征变元不是其他推理的本征变元，且除该推理的直接子推导外，不出现在~$\pi$ 的任何其他地方；否则称其为\emph{脏的}。
设 $n$ 是~$\pi$ 中脏的本征变元推理的数量。我们将通过对~$n$ 进行归纳来证明，$\pi$ 可以转化为所需的正则推导~$\pi'$。如果 $n=0$，则 $\pi$ 中所有本征变元推理都是干净的，因此 $\pi$ 是正则的，无需再证。否则，选取一个最上方的本征变元推理，假设是 \RightR{\lforall}。以该推理结尾的~$\pi$ 的子推导~$\pi_1$ 具有如下形式：
    \[
    \AxiomC{}
    \RightLabel{$\pi_2$}
    \Deduce$\Gamma \fCenter \Delta, !A(c)$
    \RightLabel{\RightR{\lforall}}
    \UnaryInf$\Gamma \fCenter \Delta, \lforall[x][!A(x)]$
    \DisplayProof
    \]
注意，由于 $c$ 是一个本征变元，它不出现在~$\Gamma$、$\Delta$ 或 $\lforall[x][!A(x)]$ 中。推导 $\pi_2$ 是正则的，因为它根本不包含任何本征变元推理。令 $a$ 为一个不在~$\pi$ 中出现的常项符号。根据 \olref{lem:subst-regular}，$\Subst{\pi_2}{a}{c}$ 是 $\Gamma \fCenter \Delta,
!A(a)$ 的一个正则推导，我们可以对其应用~$\RightR{\lforall}$。如果我们在~$\pi$ 中用推导 $\pi_1'$ 替换 $\pi_1$，
    \[
    \AxiomC{}
    \RightLabel{$\Subst{\pi_2}{a}{c}$}
    \Deduce$\Gamma \fCenter \Delta, !A(a)$
    \RightLabel{\RightR{\lforall}}
    \UnaryInf$\Gamma \fCenter \Delta, \lforall[x][!A(x)]$
    \DisplayProof
    \]
我们就将一个脏的本征变元推理替换成了一个干净的，唯一的区别是将本征变元 $c$ 替换为~$a$。所得推导具有 $n-1$ 个脏的本征变元推理，根据归纳假设即得结论。
\end{proof}

我们不会显式地引用刚才证明的命题，但将假定所有被考虑的推导都是正则的——如果它们不是，我们总可以通过替换本征变元使其正则。

\end{document}